\uuid{64Jc}
\titre{DL en l'infini/Calcul de limites/Asymptotes}
\auteur{hanani ; mzari}
\organisation{exo7}
\licence{CC BY-NC-SA 4.0}
\datecreate{2020-01-20}
\niveau{L1}
\module{Analyse}
\chapitre{Développement limité}
\sousChapitre{Applications}
\difficulte{3}
\theme{Applications des DL, DL en l'infini, Calcul de limites, Asymptotes, asymptote}
\exo7qcmid{lille-2.developpements-limites.046}

\contenu{
\texte{
Soit $a\in \Rr$ et $f$ une fonction. On dit que $f(x)$ est un $o\left((x-a)^n\right)$ au voisinage de $a$ et on écrit $f(x)=o((x-a)^n)$ si $\displaystyle \lim _{x\to a}\frac{f(x)}{(x-a)^n}=0$. On note par $DL_n(a)f(x)$ le développement limité de $f$ à l'ordre $n$ en $a$. On note aussi $DL_n(a^+)f(x)$ (resp. $DL_n(a^-)f(x)$) le développement limité de $f$ à l'ordre $n$ à droite en $a$ (resp. à gauche en $a$).
}

\begin{question}
\qclassification{125.03}
\qtags{motcle=DL en l'infini/Calcul de limites/Asymptotes}
\qtags{motcle=asymptote}
Soit $\displaystyle f(x)=x\arctan \left(\frac{1}{\sqrt{1+x^2}}\right)$. On considère la fonction $g$ telle que $\displaystyle g(x)=f\left(\frac{1}{x}\right)$ et on note $\Gamma$ le graphe de $f$. Parmi les affirmations suivantes, cocher celles qui sont vraies :
\begin{answers}
\good{$\displaystyle DL_2(0^+)g(x)=1-\frac{5x^2}{6}+o(x^2)$.}
\bad{$\displaystyle \lim _{x\to 0}g(x)=1$ et $g$ se prolonge par continuité en $0$.}
\good{$\Gamma$ admet la droite d'équation $\displaystyle y=1$ comme asymptote au voisinage de $+\infty$.}
\bad{$\Gamma$ admet la droite d'équation $\displaystyle y=1$ comme asymptote au voisinage de $-\infty$.}
\end{answers}
\vskip2mm
\begin{explanations}
On a : $\displaystyle g(x)=\frac{1}{x}\arctan \left(\frac{x}{\sqrt{1+x^2}}\right)$ si $x>0$ et $\displaystyle g(x)=\frac{-1}{x}\arctan \left(\frac{x}{\sqrt{1+x^2}}\right)$ si $x<0$. Dans le $DL_3(0)\arctan u$, on pose $\displaystyle u=\frac{x}{\sqrt{1+x^2}}=x-\frac{x^3}{2}+o(x^3)$. Ceci donne :
$$DL_2(0^+)g(x)=1-\frac{5x^2}{6}+o(x^2)\Rightarrow DL_2(+\infty)f(x)=1-\frac{5}{6x^2}+o\left(\frac{1}{x^2}\right)$$
et
$$DL_2(0^-)g(x)=-1+\frac{5x^2}{6}+o(x^2)\Rightarrow DL_2(-\infty)f(x)=-1+\frac{5}{6x^2}+o\left(\frac{1}{x^2}\right).$$
Ainsi $\Gamma$ admet la droite d'équation $\displaystyle y=1$ (resp. $y=-1$) comme asymptote au voisinage de $+\infty$ (resp. $-\infty$). Enfin, $g$ n'admet pas de limite en $0$ car $\displaystyle \lim _{x\to 0^+}g(x)=1$ et $\displaystyle \lim _{x\to 0^-}g(x)=-1$.
\end{explanations}
\end{question}
}
